% Extracto íntegro por residencia del propietario científico definitivo de masas.
\section{Refinamiento explícito en las torres \texorpdfstring{\(90/120\)}{90/120}}

La existencia del semigrupo \(\langle90,120\rangle\) no basta por sí sola:
debe construirse el mapa que asigna coeficientes a una ruta enriquecida. Sea
\[
 \mathcal H_{90,120}
 =\{90a+120b:a,b\in\mathbb N_0,\ (a,b)\ne(0,0)\}
 =\{90,120\}\cup\{30r:r\geq6\}
 =\{h_0<h_1<h_2<\cdots\}
\]
su enumeración estrictamente creciente, donde \(\mathbb N_0=\{0,1,2,\ldots\}\),
y sea \(s_{h_n}>0\) la normalización declarada del modo correspondiente. El
registro compatible de una ruta
\(\Gamma\) proporciona una sucesión de datos enteros finitos
\[
 \mathcal L(\Gamma)=(\ell_0(\Gamma),\ell_1(\Gamma),\ldots),
\]
en la que cada \(\ell_n\) reúne retorno, acarreo, devanado, memoria y frontera
en el refinamiento \(n\). Para fijar sin ambigüedad la codificación, sea
\[
 \nu(z)=
 \begin{cases}2z,&z\geq0,\\-2z-1,&z<0,\end{cases}
 \qquad
 \pi_{\rm C}(u,v)=\frac{(u+v)(u+v+1)}2+v,
\]
y defínase \(\iota:\mathbb Z^d\to\mathbb N_0\) por aplicación recursiva de
\(\pi_{\rm C}\) a \(\nu(z_1),\ldots,\nu(z_d)\). Entonces
\[
 q_n(\Gamma)
 =\frac{\iota(\ell_n(\Gamma))}
 {1+|\iota(\ell_n(\Gamma))|}\in(-1,1).
\]
La realización mínima de torre es
\[
 \eta_{h_n}(\Gamma)
 =\frac{3^{-(n+1)}q_n(\Gamma)}{s_{h_n}},
 \qquad
 \mathfrak T(\Gamma)
 =(\eta_h(\Gamma))_{h\in\mathcal H_{90,120}}.
\]

\begin{theorem}[Convergencia, naturalidad y ceguera metrológica]
Para toda ruta enriquecida,
\[
 \sum_{h\in\mathcal H_{90,120}}|\eta_h(\Gamma)s_h|
 \leq\sum_{n\ge0}3^{-(n+1)}=\frac12.
\]
Si el refinamiento prolonga el registro por prefijos, \(\mathfrak T\) es natural
respecto de las truncaciones. Además, ningún valor de masa externo interviene
en \(\mathfrak T\), y la cola después de \(N\) modos satisface
\[
 \sum_{n>N}|\eta_{h_n}s_{h_n}|
 \leq\frac{3^{-(N+1)}}2.
\]
\end{theorem}

\begin{proof}
Se tiene \(|q_n|<1\); las dos cotas son sumas geométricas. La compatibilidad
por prefijos conserva todos los coeficientes ya construidos y añade los del
nuevo nivel, lo cual da el cuadrado de naturalidad con la truncación de
\(\mathcal E_{90,120}\). La fórmula sólo contiene el registro, la codificación
explícita y la normalización modal declarada.
\end{proof}

La construcción anterior fija una realización mínima explícita y convergente,
relativa a la codificación \(\iota\) y a las normalizaciones \(s_h\) declaradas.
Una dinámica que modifique sus pesos debe declarar el operador que efectúa esa
modificación y demostrar de nuevo naturalidad y convergencia; no puede
reconstruir los pesos a partir de un residual experimental.

\section{Operadores de fibra y ley bidireccional}

Para \(F_X=\mathscr S(X)\), sea
\[
 \mathcal H_X=\ell^2(F_X)\otimes V_{\rho_X}
\]
el espacio de realización, donde \(V_{\rho_X}\) porta la representación
interna correspondiente. Cada ruta \(\gamma\in F_X\) posee firma
\[
 \sigma(\gamma)
 =(n_A,s_C,k_{\Delta_4},b_{120},b_\zeta)\in\mathbb Z^5
\]
y coeficientes de refinamiento
\(\eta(\gamma)\in\mathcal E_{90,120}\). El operador basal es
\[
 \widehat M_X^{(0)}
 =\sum_{\gamma\in F_X}
 m_{\rm clk}\,
 \chi_{\rm full}(\sigma(\gamma),\eta(\gamma))
 |\gamma\rangle\langle\gamma|\otimes I_{V_{\rho_X}}.
\]
La escala \(m_{\rm clk}\) desciende de la rama absoluta
acción--reloj--masa, incluida la década \(10^{-34}\); no se extrae de una
masa experimental.

Los lectores \(K_D\) y \(K_\Omega\) inducen operadores diagonales exactos
\(B_X^D\) y \(B_X^\Omega\). Las dos orientaciones publican
\[
 \widehat M_X^{\beta,r}
 =R_{\rm act}^{B_X^r/2}\widehat M_X^{(0)}
  R_{\rm act}^{B_X^r/2},
 \qquad r\in\{D,\Omega\}.
\]
Ambas son positivas y conmutan en la base de rutas. Por el teorema espectral
conjunto existe una única medida espectral proyectiva
\(E_X^{D,\Omega}\) sobre \(\mathbb R_+^2\) tal que
\[
 \widehat M_X^{\beta,D}
 =\int x\,dE_X^{D,\Omega}(x,y),
 \qquad
 \widehat M_X^{\beta,\Omega}
 =\int y\,dE_X^{D,\Omega}(x,y).
\]
En la base de rutas,
\[
 E_X^{D,\Omega}(A)|\gamma\rangle
 =\mathbf 1_A\!\left(m_D(\gamma),m_\Omega(\gamma)\right)
  |\gamma\rangle.
\]
Ésta es la salida bidireccional canónica: conserva el par completo y no
supone de antemano un funcional escalar.

Sea
\[
 \boldsymbol\beta_X(\gamma)
 =(m_D(\gamma),m_\Omega(\gamma)).
\]
La medida imagen uniforme de la fibra es
\[
 \nu_X=(\boldsymbol\beta_X)_\#
 \left(\frac1{|F_X|}\sum_{\gamma\in F_X}\delta_\gamma\right)
 =\frac1{|F_X|}\sum_{\gamma\in F_X}
 \delta_{(m_D(\gamma),m_\Omega(\gamma))}.
\]
Su transformada y sus momentos mixtos son
\[
 Z_X(s,t)=\int e^{sx+ty}\,d\nu_X(x,y),
 \qquad
 \mathfrak m_{p,q}(X)=\int x^py^q\,d\nu_X(x,y).
\]
Como \(\nu_X\) tiene soporte finito, \(Z_X\), o equivalentemente la familia
completa de momentos mixtos, determina el multiconjunto conjunto de pares
lectores hasta permutación de rutas. La medida imagen es, por tanto, un
invariante completo del par operatorial; no es una media de masas.

Para un canal \(P_{X,a}\), la medida espectral comprimida
\[
 F_{X,a}^{D,\Omega}(A)
 =\left.P_{X,a}E_X^{D,\Omega}(A)P_{X,a}\right|_{P_{X,a}\mathcal H_X}
\]
es una POVM sobre el subespacio preparado. Para una densidad \(\rho_X\)
soportada en ese subespacio, la distribución escalar conjunta es
\[
 \mu_{X,a}^{D,\Omega}(A)
 =\operatorname{Tr}\!\left(
 \rho_X F_{X,a}^{D,\Omega}(A)
 \right).
\]
Todo funcional declarado \(f:\mathbb R_+^2\to\mathbb R_+\) induce
\(\mu_{X,a}^f=f_\#\mu_{X,a}^{D,\Omega}\) sin modificar la PVM de origen.

Si \(u:F_X\to F_Y\) es una biyección que conserva registro, orientación y
ambos lectores, el unitario \(U_u|\gamma\rangle=|u\gamma\rangle\) satisface
\[
 U_uE_X^{D,\Omega}(A)=E_Y^{D,\Omega}(A)U_u,
 \qquad
 \nu_X=\nu_Y.
\]
La misma igualdad vale para las distribuciones físicas cuando densidad y
proyector se transportan por \(U_u\). Esta naturalidad se afirma sólo para
los entrelazadores efectivamente construidos; una coincidencia de
cardinalidades no crea por sí sola una equivalencia de fibras.

Cuando la representación física exige simetría bajo el intercambio
\(D\leftrightarrow\Omega\), el agregado operatorial mínimo es la media
geométrica
\[
 \widehat M_X^{\rm sym}
 =\widehat M_X^{\beta,D}\#\widehat M_X^{\beta,\Omega},
 \qquad
 A\#B
 =A^{1/2}\bigl(A^{-1/2}BA^{-1/2}\bigr)^{1/2}A^{1/2},
\]
sobre el soporte positivo común. En el caso diagonal,
\[
 \widehat M_X^{\rm sym}
 =R_{\rm act}^{(B_X^D+B_X^\Omega)/4}
  \widehat M_X^{(0)}
  R_{\rm act}^{(B_X^D+B_X^\Omega)/4}.
\]
Esta fórmula es consecuencia de la hipótesis explícita de intercambio; no
sustituye la medida conjunta cuando tal simetría no ha sido establecida.

\begin{theorem}[Caracterización del agregado simétrico]
La media geométrica es positiva, simétrica, covariante por congruencia para
toda \(S\) invertible y autodual:
\[
 A\#B=B\#A,
 \qquad
 S^*(A\#B)S=(S^*AS)\#(S^*BS),
 \qquad
 (A\#B)^{-1}=A^{-1}\#B^{-1}.
\]
Además, \(A\#A=A\). Por consiguiente, en los sectores donde existe un
entrelazador físico que intercambia los lectores, el agregado no introduce un
parámetro ni privilegia una orientación.
\end{theorem}

\begin{proof}
La identidad de congruencia se entiende para \(S\) invertible. Las identidades
resultan del cálculo funcional de la raíz positiva de
\(A^{-1/2}BA^{-1/2}\). En el caso conmutativo se reducen a
\(A\#B=(AB)^{1/2}\).
\end{proof}

Existe además un agregado escalar canónico bajo axiomas distintos. Si
\(s_n\) es continuo, invariante por permutación, multiplicativo coordenada a
coordenada y normalizado por \(s_n(cI)=c\), entonces
\[
 s_n(T)=\det(T)^{1/n}.
\]
En efecto, al pasar a logaritmos, continuidad y aditividad hacen lineal el
funcional sobre los logaritmos de los autovalores; la simetría iguala todos
los coeficientes y la normalización fija su valor en \(1/n\). Para el par de
lectores, la simetría de intercambio da el invariante
\[
 S_{D,\Omega}
 =\sqrt{
 \operatorname{ndet}(T_D)\operatorname{ndet}(T_\Omega)},
 \qquad
 \operatorname{ndet}(T)=\det(T)^{1/n}.
\]
Este resultado clasifica la fibra completa bajo los axiomas enunciados. No se
identifica automáticamente con una línea de masa: dicha identificación exige
que el acoplamiento físico adopte precisamente ese funcional.

Para la fibra electrónica de rango uno, la medida conjunta es puntual, los dos
lectores comparten el primer componente \(80\), la compresión es escalar y se
recupera el cierre beta
\[
 m_e^\beta
 =0.5109989506900008086082571648\ldots\ {\rm MeV}.
\]
El valor basal permanece como etapa de construcción; no reemplaza al cierre
beta.

\section{Escalarización por simetría y acoplamiento}

Sea \(G_X\) el grupo compacto de simetría no rota de la fibra y
\(U_X:G_X\to\mathcal U(\mathcal H_X)\) su representación. La esperanza
condicional de simetría se define por
\[
 \mathbb E_X(T)
 =\int_{G_X}U_X(g)TU_X(g)^*\,d\mu_X(g),
\]
donde \(\mu_X\) es la medida de Haar normalizada. El operador físico de masa
antes de abrir canales se define, cuando el tipo de estado declara un
funcional \(\Phi_X:\mathbb R_+^2\to\mathbb R_+\), por
\[
 \mathcal M_X
 =\mathbb E_X\!\left(
   \Phi_X(\widehat M_X^{\beta,D},\widehat M_X^{\beta,\Omega})
  \right).
\]
El cálculo funcional se realiza con la medida conjunta
\(E_X^{D,\Omega}\). En una representación que contiene un entrelazador de
intercambio de lectores puede tomarse \(\Phi_X(x,y)=\sqrt{xy}\). Si el tipo
físico no declara \(\Phi_X\), la salida completa sigue siendo la PVM conjunta;
no se inventa un escalar.

\begin{theorem}[Cierre escalar, multiplete y distribución]
El operador \(\mathcal M_X\) pertenece al conmutante
\(U_X(G_X)'\). Si
\[
 \mathcal H_X\simeq
 \bigoplus_\lambda V_\lambda\otimes\mathbb C^{m_\lambda},
\]
entonces
\[
 \mathcal M_X=
 \bigoplus_\lambda I_{V_\lambda}\otimes M_{X,\lambda}.
\]
En un sector irreducible de multiplicidad \(1\), la restricción es un escalar.
En un sector reducible, la salida exacta es el espectro de los bloques
\(M_{X,\lambda}\), con sus multiplicidades. No existe una indeterminación:
existe un codominio distinto.
\end{theorem}

\begin{proof}
La invariancia de la medida de Haar implica
\(U_X(h)\mathbb E_X(T)U_X(h)^*=\mathbb E_X(T)\) para todo \(h\in G_X\).
La descomposición sigue del teorema de completa reducibilidad y del lema de
Schur. El caso irreducible produce un múltiplo de la identidad; el caso
reducible conserva exactamente los bloques de multiplicidad.
\end{proof}

Sea \(P_{X,a}\) el proyector determinado por el estado preparado y el canal o
aparato físico \(a\). La medida espectral observable es
\[
 \mu_{X,a}(B)=
 \langle\Psi_X,
 P_{X,a}E_{\mathcal M_X}(B)P_{X,a}\Psi_X\rangle.
\]
Hay una línea escalar precisamente cuando el soporte espectral comprimido es
puntual:
\[
 E_{\mathcal M_X}(\{m_X\})P_{X,a}=P_{X,a},
 \qquad\text{equivalentemente}\qquad
 (\mathcal M_X-m_X I)\,P_{X,a}=0.
\]
La igualdad aislada
\(P_{X,a}\mathcal M_X P_{X,a}=m_X P_{X,a}\) sólo fija una esperanza si el
subespacio no es reductor y no basta para demostrar una línea espectral.
Si la compresión conserva varios autovalores, la salida es un multiplete o una
distribución. Sustituirla por una media numérica destruiría información
física y genealógica.


